\section{Platform and Custom Hardware}
\label{sec:platform}

The watch limits this project depends on:
\begin{itemize}
\item \textbf{C1}, memory: 512\,KB SRAM, 8\,MB PSRAM, 32\,MB flash,
vendor wiki, constraint -- bounds the size of the trained
on-device classifier.
\item \textbf{C2}, wireless: 2.4\,GHz Wi-Fi (802.11\,b/g/n), vendor
wiki, constraint -- the link this project reports over and
degrades for testing. (The product page also lists 802.11ax; the
wiki figure is used here as the more detailed of the two sources.)
\item \textbf{C3}, sensing: 6-axis IMU (accelerometer + gyroscope),
vendor wiki, constraint -- source of the watch's own wrist motion
data.
\item \textbf{C4}, power: 400\,mAh Li-ion battery, 3.7\,V nominal,
vendor wiki (battery as supplied), constraint -- bounds how often
the board and watch can transmit.
\end{itemize}

The custom hardware is a low-fidelity prototype: an off-the-shelf
microcontroller board with its own Wi-Fi radio, carrying one or two
tethered inertial sensor modules (chest and/or thigh), wired to the
board over I2C. The board does what the watch cannot: sense motion
away from the wrist. It connects to the watch wirelessly, the same
way it connects to the companion computer, rather than through the
watch's I2C/UART/USB expansion pads, which this design does not use.
This wireless link is the interface this project depends on, and
is stated as a numbered requirement in Section~\ref{sec:reqs}.

The board extracts motion features (for example, magnitude and
cadence) from its chest and/or thigh sensors and sends them to the
watch. The watch extracts features from its own wrist IMU, aligns
the board's features to its own by timestamp, and combines them
into one classification of sedentary, walking or running. If a
board stream is missing, the watch classifies on the features it
still has, down to wrist alone, and tags each report with which
sensors produced it.

In Project II, the breadboard/perfboard prototype is replaced by
one custom PCB: a carrier board that the microcontroller board
plugs into, adding battery charging, a power switch and connectors
for the tethered chest and/or thigh sensors. The PoC in
Section~\ref{sec:poc} is run on the low-fidelity prototype, so its
results do not depend on the PCB being ready.
